\section{Large average-free sets}
\label{sec:average-free-sets}

The hashing argument uses the following form of the no-carry sphere
construction. Generalized average-free sets of this kind occur in
\cite{VC2015,CVZ2019}. We include the construction and its size estimate
to specify exactly the property needed here.

\begin{lemma}[Average-free sets]
\label{lem:average-free}
Fix an integer $L\geq 2$. For every sufficiently large integer $q$ there
is a set
\[
 B\subseteq\left\{0,\ldots,\left\lfloor\frac{q-1}{L}\right\rfloor\right\}
\]
such that
\[
 x_1+\cdots+x_L=Ly,\qquad x_1,\ldots,x_L,y\in B,
 \quad\Longrightarrow\quad x_1=\cdots=x_L=y,
\]
and
\[
 \log|B|=\log q-O_L\!\left(\sqrt{\log q}\right).
\]
The integer $q$ need not be prime in this lemma.
\end{lemma}

\begin{proof}
Put
\[
 \ell=\left\lfloor\sqrt{\log q}\right\rfloor,
 \qquad
 m=\left\lfloor\frac{(q/(2L))^{1/\ell}}{L}\right\rfloor,
 \qquad b=Lm.
\]
For sufficiently large $q$, we have $\ell\geq 1$, $m\geq 2$, and
$b^\ell\leq q/(2L)$. There are $m^\ell$ vectors in
$\{0,\ldots,m-1\}^{\ell}$ and at most $\ell(m-1)^2+1$ possible values
of their squared Euclidean norms. Choose a norm level $\mathcal C$
with
\[
 |\mathcal C|\geq\frac{m^\ell}{\ell(m-1)^2+1},
\]
and encode its vectors in base $b$:
\[
 B=\left\{\sum_{t=0}^{\ell-1}a_t b^t:
                (a_0,\ldots,a_{\ell-1})\in\mathcal C\right\}.
\]
The encodings are distinct and lie between $0$ and $b^\ell-1$.
Since $b^\ell\leq q/(2L)$, they belong to the required interval.

Suppose that $x_1+\cdots+x_L=Ly$ for elements of $B$, and let
$a^{(1)},\ldots,a^{(L)},c\in\mathcal C$ be their digit vectors.
Every digit sum on either side is at most $L(m-1)<b$. Thus neither
side has carries, and uniqueness of the base-$b$ expansion gives
\[
 a^{(1)}+\cdots+a^{(L)}=Lc.
\]
All these vectors have the same squared norm. Consequently
\[
 \frac1L\sum_{i=1}^L\lVert a^{(i)}-c\rVert^2
 =\frac1L\sum_{i=1}^L\lVert a^{(i)}\rVert^2-\lVert c\rVert^2
 =0.
\]
Each $a^{(i)}$ equals $c$, proving the average-free property.

For the size estimate, write $t=\log q$. Then $\ell\asymp\sqrt t$
and
\[
 \log m=\frac{t-\log(2L)}{\ell}-\log L+O_L(m^{-1})
       =\sqrt t+O_L(1).
\]
In particular $m\to\infty$, and
\[
 \ell\log m=t-\log(2L)-\ell\log L+O_L(\ell/m)
             =t-O_L(\sqrt t).
\]
Moreover,
\[
 \log\bigl(\ell(m-1)^2+1\bigr)
 \leq \log(\ell+1)+2\log m=O_L(\sqrt t).
\]
The lower bound for $|\mathcal C|=|B|$ therefore gives
$\log|B|\geq\log q-O_L(\sqrt{\log q})$. The opposite estimate
$|B|\leq q$ completes the proof.
\end{proof}

\section{Hashes and exact survival probabilities}
\label{sec:hashing}

We retain the notation $S\subseteq\mathbb Z^k$, $M$, $h$, $E_j$,
and $\beta$ from Section~\ref{sec:fibers}, with $k\geq 3$ fixed. In
particular, every column in $S$ has coordinate sum zero. We use the
tight-support hashing construction of \cite{CVZ2019}; the uniform rank
bound and the exact probabilities needed here are established below.

\begin{lemma}[Uniform reduction of ranks]
\label{lem:rank-reduction}
Let
\[
 C=\max_{s,s'\in S}\max_{1\leq i\leq k}|s_i-s'_i|,
 \qquad
 B_0>\max\{k,\,k!\max\{1,C\}^{k}\}.
\]
For every prime $q>B_0$, any integer matrix whose columns belong to
$S-S$ has the same rank over $\Q$ and over $\F_q$. The threshold is
independent of the number of columns.
\end{lemma}

\begin{proof}
If the rational rank is zero, every entry is zero and the assertion
is immediate. Otherwise, let that rank be $r$. There is a nonzero
integer $r\times r$ minor, of absolute value at most
\[
 r!C^r\leq k!\max\{1,C\}^{k}<q.
\]
It remains nonzero modulo $q$, so the rank over $\F_q$ is at least
$r$. Every larger minor is zero as an integer and remains zero
modulo $q$, proving the opposite inequality. The argument depends
only on the entries and the number $k$ of rows.
\end{proof}

Fix a prime $q>B_0$ large enough for Lemma~\ref{lem:average-free}
with $L=k-1$, and fix a positive multiple $N$ of $M$. Choose the set
$B$ from that lemma, viewed also as a subset of $\F_q$ by reduction
modulo $q$. Choose
\[
 v=(v_1,\ldots,v_N)\in\F_q^N,
 \qquad u_1,\ldots,u_{k-1}\in\F_q
\]
independently and uniformly. For an integer-labelled local word
$w=(w_1,\ldots,w_N)$, define
\begin{align*}
 b_i(w)&=u_i+\sum_{t=1}^Nw_t v_t &&(1\leq i<k),\\
 b_k(w)&=\frac{\sum_{i<k}u_i-\sum_{t=1}^Nw_t v_t}{k-1}.
\end{align*}
These expressions are in $\F_q$; the denominator is invertible
because $q>k$.

An original support edge $I\in S^N$ may be described by its columns
$I(t)\in S$ or by its $k$ local words $I_1,\ldots,I_k$. The
zero-sum condition on every column gives
\begin{equation}
 \sum_{i<k}b_i(I_i)=(k-1)b_k(I_k)\qquad\text{in }\F_q.
 \label{eq:tight-hash-relation}
\end{equation}
Keep in each mode exactly the local basis words whose hash belongs
to $B$. This is a coordinate restriction of all $k$ factors. If
$I$ survives, represent each of its hashes by the corresponding
integer in $B$. Both sides of \eqref{eq:tight-hash-relation} then
belong to $[0,q-1]$, so their congruence is an equality of integers.
Lemma~\ref{lem:average-free} implies that all $k$ hashes equal a
single element of $B$. Conversely, this condition implies survival.

For a fixed edge $I$ and a prescribed common hash $b\in B$, each
choice of $v$ permits exactly one offset tuple: set
$u_i=b-\sum_t I_i(t)v_t$ for $i<k$. The last hash then also equals
$b$ by \eqref{eq:tight-hash-relation}. The events for distinct $b$
are disjoint. Of the $q^{N+k-1}$ choices of all hashing variables,
exactly $|B|q^N$ make $I$ survive. Thus
\begin{equation}
 \Prob[I\text{ survives}]=|B|q^{-(k-1)}.
 \label{eq:single-survival}
\end{equation}

Now let distinct original edges $I,J\in S^N$ share a local word,
say $I_j=J_j$. Put
\[
 V=\operatorname{span}_{\Q}
       \{I(t)-J(t):1\leq t\leq N\},
 \qquad r=\dim V.
\]
Every column difference has coordinate sum zero and coordinate $j$
equal to zero. Hence $V\subseteq E_j$ and $r\leq k-2$. Also
$r\geq 1$, since the integer labelling is injective and $I\ne J$.

If both edges survive, their common hash values agree: the shared
local word has the same hash for both. For a prescribed common
$b\in B$, the offsets are uniquely determined by $I$ once $v$ is
chosen. The additional conditions on $J$ are precisely
\[
 \sum_{t=1}^N\bigl(I_i(t)-J_i(t)\bigr)v_t=0,
 \qquad 1\leq i<k.
\]
The rows of the full $k\times N$ difference matrix sum to zero.
Its last row is therefore redundant, over both $\Q$ and $\F_q$.
Deleting that row leaves rank $r$ by
Lemma~\ref{lem:rank-reduction}. The displayed system has exactly
$q^{N-r}$ solutions for $v$, each with one associated offset tuple.
Consequently
\begin{equation}
 \Prob[I\text{ and }J\text{ survive}]
       =|B|q^{-(k-1+r)}.
 \label{eq:pair-survival}
\end{equation}
This argument also applies when the shared mode is $j=k$: its
shared hash still forces equality of the two common values, and
the last row remains redundant. Neither edge is required to have
uniform joint type for \eqref{eq:single-survival} or
\eqref{eq:pair-survival}.

\section{Isolation against the full original support}
\label{sec:isolation}

We now distinguish the counting class $\mathcal A_N$ from the
support $S^N$ of the full tensor power. Let $X$ count surviving
candidates $I\in\mathcal A_N$. Let $Y$ count ordered pairs $(I,J)$
such that $I$ is a surviving candidate, $J\ne I$ is any surviving
original edge, and $I$ and $J$ share a local word in at least one
mode. In particular, no type condition is imposed on $J$.
Equation~\eqref{eq:single-survival} gives
\begin{equation}
 \E X=|\mathcal A_N|\,|B|q^{-(k-1)}.
 \label{eq:candidate-expectation}
\end{equation}

Let $\mathcal V$ be the collection of distinct nonzero rational
spans of subsets of $S-S$ that are contained in at least one $E_j$,
and put $K=|\mathcal V|$. This collection is finite because $S-S$
is finite; it depends on $\mathbf d$, but not on $N$. Every
$V\in\mathcal V$ has dimension between $1$ and $k-2$, and each
pair counted by $Y$ has exactly one such difference span.

Fix $I\in\mathcal A_N$ and $V\in\mathcal V$ of dimension $r$.
The number of partners $J\in S^N$ whose column differences with
$I$ span exactly $V$ is at most the number whose differences all
belong to $V$. By \eqref{eq:fiber-count} and
Lemma~\ref{lem:fiber-entropy}, this latter number is at most
$\exp(Nr\beta)$. We use a mode $j$ with $V\subseteq E_j$ only to
apply that entropy estimate; we do not sum over choices of $j$.
Summing \eqref{eq:pair-survival} over the candidates and their
possible difference spans therefore yields
\begin{equation}
 \E Y\leq \E X
       \sum_{V\in\mathcal V}
          \exp\bigl(N\dim(V)\beta\bigr)q^{-\dim(V)}.
 \label{eq:collision-expectation}
\end{equation}
The larger fiber count is an upper bound even though it also
allows partners of smaller difference span. Each actual pair is
assigned to its unique exact span in this sum, so there is no
extra multiplicity for the shared mode.

For a realization of the hashing variables, call a surviving
candidate \emph{isolated} if it shares no local word with any other
surviving original edge. Denote the set of these candidates by
$\mathcal I$. Every nonisolated candidate contributes at least one
ordered pair to $Y$, and hence
\[
 |\mathcal I|\geq X-Y.
\]
This is the candidate-indicator instance of the weighted
global-isolation principle from \cite{VC2019}. The point needed
for a coordinate restriction is that the comparison uses the
entire original support, as follows.

If $\mathcal I$ is empty, then $X-Y\leq 0$ and the corresponding
lower estimate follows already from $\Qmon\geq 1$. Otherwise,
retain exactly the local words belonging to the edges in
$\mathcal I$. Distinct edges of $\mathcal I$ have distinct local
words in every mode, since they are isolated. Any original edge
remaining after this final coordinate restriction would have
survived the earlier hash restriction. It also shares a retained
word with some $I\in\mathcal I$. If it differed from $I$, it would
contradict the definition of isolation. Thus the remaining support
is exactly $\mathcal I$, and independent relabelling in the $k$
modes gives the diagonal tensor of size $|\mathcal I|$.

The hashing probability space is finite, so some realization has
$X-Y\geq\E(X-Y)$. We have proved
\begin{equation}
 \Qmon\bigl(D_{\mathbf d}^{\boxtimes N}\bigr)\geq\E(X-Y).
 \label{eq:expected-diagonal}
\end{equation}
No local restriction selecting $\mathcal A_N$ was used. Isolation
only within the candidate class would not justify the conclusion:
an edge of $S^N\setminus\mathcal A_N$ could remain on the retained
local words.

\section{The entropy lower bound}
\label{sec:lower-bound}

Fix $\delta>0$ and set $\lambda=\beta+\delta>0$. Choose primes
$q$ tending to infinity, and for sufficiently large $q$ define
\[
 N=N(q)=M\left\lfloor\frac{\log q}{M\lambda}\right\rfloor.
\]
Then $N$ is a positive multiple of $M$, tends to infinity with
$q$, and satisfies
\begin{equation}
 N\lambda\leq\log q<(N+M)\lambda.
 \label{eq:prime-first-subsequence}
\end{equation}
The primes are chosen first and the tensor powers second. Thus
unboundedness of the primes suffices; no estimate for prime gaps
is needed. For all sufficiently large $q$, both the rank threshold
and Lemma~\ref{lem:average-free} apply.

For $r=\dim V\geq 1$, equation~\eqref{eq:prime-first-subsequence}
gives
\[
 \exp(Nr\beta)q^{-r}\leq\exp(-Nr\delta)
                       \leq\exp(-N\delta).
\]
It follows from \eqref{eq:collision-expectation} that
\[
 \E Y\leq\E X\,K\exp(-N\delta).
\]
Combining this inequality with \eqref{eq:candidate-expectation}
and \eqref{eq:expected-diagonal} yields
\begin{equation}
 \Qmon\bigl(D_{\mathbf d}^{\boxtimes N}\bigr)
 \geq\bigl(1-Ke^{-N\delta}\bigr)
        |\mathcal A_N|\,|B|q^{-(k-1)}
 \label{eq:finite-monomial-bound}
\end{equation}
for these powers and all sufficiently large $q$. The prefactor
is positive and tends to one.

All of $\mathbf d,k,M,K,\beta,\delta$ are fixed before $q$ tends
to infinity. By \eqref{eq:type-count},
$\log|\mathcal A_N|=Nh-o(N)$. Lemma~\ref{lem:average-free} gives
$\log|B|=\log q-O_k(\sqrt{\log q})$, and
\eqref{eq:prime-first-subsequence} gives
$\log q=N(\beta+\delta)+O_{\mathbf d,\delta}(1)$.
In particular $\log q=O_{\mathbf d,\delta}(N)$, so the error in
the average-free estimate is $o(N)$. The logarithm of the right
side of \eqref{eq:finite-monomial-bound} is consequently
\begin{align*}
 Nh-(k-2)\log q-o(N)
 &=N\bigl[h-(k-2)(\beta+\delta)\bigr]-o(N)\\
 &=N\bigl[\HH(p)-(k-2)\delta\bigr]-o(N).
\end{align*}
Here we used $(k-2)\beta=h-\HH(p)$.

Lemma~\ref{lem:limits} transfers this lower bound along the
unbounded sequence of powers to the asymptotic limit, giving
\[
 \log\Qmt(D_{\mathbf d})\geq\HH(p)-(k-2)\delta.
\]
Since $\delta>0$ was arbitrary, and every coordinate restriction
is a linear restriction, we conclude that
\begin{equation}
 \Qt(D_{\mathbf d})\geq\Qmt(D_{\mathbf d})
                    \geq\exp\bigl(\HH(p)\bigr).
 \label{eq:monomial-lower}
\end{equation}
